\section{Proofs of propositions in \texorpdfstring{\Cref{sec:perturbation}}{Section~4}}
\label{sec:appendix-proofs}

Throughout this appendix, notation and standing assumptions
\ref{SA:primitive}--\ref{SA:non-coboundary} from \Cref{sec:perturbation}
are in force. We write \(B_\theta:=I-A_\theta/\lambda(\theta)\) and
\(R_\theta:=B_\theta^{\#}\) for the group inverse.

\begin{proof}[Proof of \Cref{thm:finite-part-reduced-determinant-group-inverse-gradient}]
  Since \(\lambda(\theta)\) is a simple eigenvalue
  \cite[Chapter~II, \S1]{Kato1995Perturbation}, \(B_\theta\) has a
  simple zero eigenvalue and all other eigenvalues are non-zero; hence
  the group inverse \(R_\theta\) exists; see
  \cite[Chapter~4]{CampbellMeyer1979}. Because the spectrum depends
  analytically, \(P_\theta\) and \(R_\theta=B_\theta^{\#}\) vary
  analytically on \(\theta\). On \((I-P_\theta)\CC^d\), the
  restriction of \(B_\theta\) is invertible and its ordinary
  determinant equals \(\det\nolimits'(B_\theta)=C(A_\theta,\lambda(\theta))^{-1}\).
  Jacobi's formula on this complementary spectral subspace gives
  \[
    \partial_\theta\log\det\nolimits'(B_\theta)
    =\Tr(R_\theta\,\partial_\theta B_\theta).
  \]
  Since \(C(A_\theta,\lambda(\theta))=\det\nolimits'(B_\theta)^{-1}\),
  \[
    \partial_\theta\log C(A_\theta,\lambda(\theta))
    =-\Tr(R_\theta\,\partial_\theta B_\theta),
  \]
  while
  \[
    \partial_\theta B_\theta
    =-\frac{\partial_\theta A_\theta}{\lambda(\theta)}
      +\frac{\partial_\theta\lambda(\theta)}{\lambda(\theta)^2}
       A_\theta.
  \]
  Substituting yields the stated formula.
\end{proof}

\begin{proof}[Proof of \Cref{prop:secondary-fredholm-normal-convergence}]
  Choose a closed disc \(U\) around \(\theta_0\) so small that
  \(\lambda(\theta)\) stays simple on \(U\) and \(q_U\in(q,1)\) is a
  uniform spectral-gap bound. Set
  \[
    \lambda_\ast:=\inf_{\theta\in U}\abs{\lambda(\theta)}>0,
    \qquad
    M_0:=\sup_{\theta\in U}\norm{A_\theta}<\infty,
  \]
  and
  \[
    M_1:=\sup_{\theta\in U}
    \bigl(\norm{\partial_\theta A_\theta}
      +\abs{\partial_\theta\lambda(\theta)}\bigr)<\infty.
  \]
  For each \(k\ge 2\), the operator \(G_k(\theta):=I-\lambda(\theta)^{-k}A_\theta\) acts by the scalar
  \(1-\lambda(\theta)^{1-k}\) on the eigenline for \(\lambda(\theta)\)
  and by \(I-\lambda(\theta)^{-k}A_\theta\) on
  \((I-P_\theta)\CC^d\). On the eigenline the scalar
  \(1-\lambda(\theta)^{1-k}\) is bounded away from zero for \(k\ge 2\).
  On the complementary subspace, the spectral radius satisfies
  \(\rho(A_\theta|_{(I-P_\theta)\CC^d})\le q_U\abs{\lambda(\theta)}\),
  so \(\lambda(\theta)^{-k}A_\theta\) has spectral radius at most
  \(q_U\abs{\lambda(\theta)}^{1-k}\to 0\) as \(k\to\infty\).
  Hence every \(G_k(\theta)\) is invertible on \(U\) for all
  sufficiently large \(k\), and for the finitely many small \(k\),
  invertibility is guaranteed by compactness.

  Define the holomorphic branch
  \[
    \Phi_k(\theta)
    :=\Phi_k(\theta_0)
      -\int_{\theta_0}^{\theta}
        \Tr\!\bigl(G_k(\xi)^{-1}\partial_\xi G_k(\xi)\bigr)\,\dd\xi,
  \]
  where \(\Phi_k(\theta_0):=\log\det G_k(\theta_0)^{-1}\) is
  computed via the principal branch and the integral is along any path
  in the simply connected \(U\).

  For \(k\) large enough that
  \(\sup_{\theta\in U}\norm{I-G_k(\theta)}\le 1/2\), the Taylor series
  \(\Phi_k(\theta)=\sum_{m\ge 1}\Tr((\lambda(\theta)^{-k}A_\theta)^m)/m\)
  converges absolutely, giving
  \[
    \abs{\Phi_k(\theta)}
    \le d\sum_{m\ge 1}\frac{(M_0/\lambda_\ast^{k})^m}{m}
    \le C\,q_U^k.
  \]
  Jacobi's formula yields
  \[
    \partial_\theta\Phi_k(\theta)
    =-\Tr\!\bigl(G_k(\theta)^{-1}\partial_\theta G_k(\theta)\bigr),
  \]
  with
  \[
    \sup_{\theta\in U}\norm{\partial_\theta G_k(\theta)}
    \le M_1/\lambda_\ast^k+kM_0M_1/\lambda_\ast^{k+1}
    \le C'\,k\,q_U^k.
  \]
  Combining gives \(\abs{\partial_\theta\Phi_k(\theta)}\le C''kq_U^k\).
  Absorbing finitely many small \(k\) values into the constant proves
  the stated bound. Normal convergence and termwise differentiation of
  \(\sum_{k\ge 2}\mu(k)\Phi_k/k\) follow from
  \(\sum_{k\ge 2}kq_U^k<\infty\).
\end{proof}

\begin{proof}[Proof of \Cref{thm:logM-holomorphic-variation}]
  For each \(\theta\) near \(\theta_0\),
  \Cref{def:algebraic-finite-part} gives
  \[
    F(A_\theta,\lambda(\theta))
    =\log C(A_\theta,\lambda(\theta))
      +\sum_{k\ge 2}\frac{\mu(k)}{k}\Phi_k(\theta).
  \]
  By \Cref{prop:secondary-fredholm-normal-convergence}, the correction
  series converges normally and may be differentiated termwise. The
  derivative of \(\log C(A_\theta,\lambda(\theta))\) is given by
  \Cref{thm:finite-part-reduced-determinant-group-inverse-gradient}.
\end{proof}

\begin{proof}[Proof of \Cref{prop:logM-truncation}]
  Let \(U\), \(C_0\), and \(q_U\) be given by
  \Cref{prop:secondary-fredholm-normal-convergence}. Then
  \[
    \sum_{k>N}\left|\frac{\mu(k)}{k}\Phi_k(\theta)\right|
    \le C_0\sum_{k>N}q_U^k,
  \]
  and the same estimate holds with \(\Phi_k\) replaced by
  \(\partial_\theta\Phi_k\). Summing the geometric tail gives the
  result.
\end{proof}

\begin{proof}[Proof of \Cref{prop:twisted-determinant-orbit-moments}]
  For each \(n\ge 1\),
  \[
    \Tr(M_t^n)=\sum_{x\in\Fix(\sigma_A^n)}e^{tS_nf(x)}.
  \]
  Applying \Cref{thm:fredholm-witt-bridge} to the finite matrix \(M_t\)
  gives the determinant identity. The primitive-orbit product is the
  corresponding Euler-product expansion by M\"obius inversion, exactly as
  in \Cref{prop:prelim-trace-primitive}. Differentiating the logarithm
  with respect to \(t\) at \(t=0\) yields the moment identities.
\end{proof}

\begin{proof}[Proof of \Cref{thm:twisted-determinant-rigidity}]
  Equality of determinants near \((0,0)\) implies equality of the
  reciprocal determinants and therefore, by
  \Cref{prop:twisted-determinant-orbit-moments}, equality of the trace
  functions
  \[
    \Tr\!\bigl((M_t^{(f)})^n\bigr)
    =\Tr\!\bigl((M_t^{(g)})^n\bigr)
    \qquad (n\ge 1)
  \]
  for \(t\) near \(0\). Differentiating \(r\) times at \(t=0\) gives
  agreement of all periodic-orbit moments.

  For the eigenvalue branch statement, let
  \(\lambda_0=\rho(M_0^{(f)})=\rho(M_0^{(g)})\) be the common simple
  Perron eigenvalue at \(t=0\). By continuity of eigenvalues and the
  spectral-rigidity statement
  \Cref{thm:fredholm-spectral-rigidity}, for each \(t\) near \(0\) the
  full spectral multisets of \(M_t^{(f)}\) and \(M_t^{(g)}\) agree.
  In particular, the unique analytic dominant eigenvalue branches
  coincide: \(\lambda_f(t)=\lambda_g(t)\).
\end{proof}

\begin{proof}[Proof of \Cref{cor:twisted-determinant-cumulants}]
  By \Cref{thm:twisted-determinant-rigidity}, the dominant
  eigenvalue branches agree near \(t=0\). Since the centered cumulant
  generating function is
  \(t\mapsto\log(\lambda_f(t)/\lambda_f(0))-t\lambda_f'(0)/\lambda_f(0)\),
  the asymptotic cumulants are the same for \(f\) and \(g\).
\end{proof}

\begin{proof}[Proof of \Cref{prop:off-origin-decay}]
  The entries of \(\widetilde M_u\) are continuous in \(u\), so the
  spectral radius \(u\mapsto\rho(\widetilde M_u)\) is continuous on
  \(\RR\). By \Cref{def:arithmetic-lattice} and
  \cite[Theorem~1]{Parry1992Cocycle},
  \(\rho(\widetilde M_u)=\lambda\) if and only if \(u\in\Lambda_f\).
  Since \(K\Subset\RR\setminus\Lambda_f\) is compact and
  \(\rho(\widetilde M_u)<\lambda\) for all \(u\in K\), the supremum
  \(\kappa_K<1\).

  Fix \(\vartheta_K\in((1+\kappa_K)/2,1)\) and set
  \(r_K:=\vartheta_K\lambda\). The holomorphic functional calculus gives
  \[
    \widetilde M_u^n
    =\frac{1}{2\pi i}\int_{\abs{z}=r_K}
      z^n(zI-\widetilde M_u)^{-1}\,\dd z,
  \]
  so
  \(\sup_{u\in K}\|\widetilde M_u^n\|\le 2\pi C_K^{(0)}r_K^n\),
  where \(C_K^{(0)}:=\sup_{u\in K,\abs{z}=r_K}\|(zI-\widetilde M_u)^{-1}\|<\infty\).
  Multiplying by \(\lambda^{-n}\|\ell_0\|\,\|r_0\|\) gives the stated
  estimate.
\end{proof}

\begin{proof}[Proof of \Cref{thm:clt-spectral}]
  The map \(u\mapsto \widetilde M_u\) is entire. Since \(B=\widetilde
  M_0\) is primitive, analytic perturbation theory gives
  an analytic simple eigenvalue \(u\mapsto\lambda_c(u)\),
  an analytic rank-one projector \(u\mapsto P_u\), and constants
  \(C>0\), \(\vartheta\in(0,1)\) such that
  \[
    \widetilde M_u^n=\lambda_c(u)^nP_u+N_u^n,
    \qquad
    \|N_u^n\|\le C(\vartheta\lambda)^n
  \]
  for \(|u|<\delta\) and \(n\ge 1\); see
  \cite[Chapter~II, \S2]{Kato1995Perturbation}. Because \(\bar f\) is
  centered, \(\lambda_c'(0)=0\). The standard spectral method gives
  \[
    \log\frac{\lambda_c(u)}{\lambda}
    =-\frac{\sigma_f^2}{2}u^2+O(u^3),
  \]
  where \(\sigma_f^2>0\) whenever \(f\) is not a coboundary modulo
  constants; see
  \cite[Chapter~5]{Ruelle1978Thermodynamic},
  \cite[Chapter~3]{Baladi2000PositiveTransfer},
  and \cite[Chapter~IV]{HennionHerve2001}. For each fixed
  \(u_0\in\RR\),
  \[
    \phi_n(u_0/\sqrt n)
    =\left(\frac{\lambda_c(u_0/\sqrt n)}{\lambda}\right)^n
      \ell_0^{\mathsf T}P_{u_0/\sqrt n}r_0
      +O(\vartheta^n)
    \longrightarrow e^{-\sigma_f^2u_0^2/2}.
  \]
  L\'evy's continuity theorem yields the central limit theorem.
\end{proof}
